\section{Experiments}
\label{sec:experiments}

This section describes the datasets (Section~\ref{sec:exp_datasets}),
baselines (Section~\ref{sec:exp_baselines}), metrics
(Section~\ref{sec:exp_metrics}), and implementation details
(Section~\ref{sec:exp_impl}) used to evaluate CSM-SAM. All tables of
numerical results live in Section~\ref{sec:results}; this section fixes
the experimental protocol that those tables instantiate. Unless
otherwise stated, every baseline and every CSM-SAM variant is evaluated
with the same preprocessing, the same train/val/test splits, and the
same metric implementations, so that differences reported in
Section~\ref{sec:results} can be attributed to the architectural axes
laid out in Section~\ref{sec:method}.

\subsection{Datasets}
\label{sec:exp_datasets}

We evaluate CSM-SAM on eight datasets organised into three groups: a
primary longitudinal medical benchmark, a secondary longitudinal
medical benchmark, two single-timepoint medical pretraining corpora,
and four bitemporal change-detection transfer benchmarks. Only the
first group carries the central aggDSC claim; the remaining datasets
probe the generality of the frozen-backbone plus cross-session-memory
recipe.

\subsubsection{Primary longitudinal benchmark}
\label{sec:exp_hnts}

\textbf{HNTS-MRG 2024 Task 2.} The primary benchmark is the
HNTS-MRG 2024 head and neck tumor segmentation
challenge~\cite{wahid2024hntsmrg}. It provides 150 training patients
with paired pre-RT (week 0) and mid-RT (week 2--3) T2-weighted MR
volumes, manual GTVp (primary tumor) and GTVn (nodal) contours, and an
organizer-provided deformably registered pre-RT mask on the mid-RT
grid. Task 2 evaluates mid-RT segmentation with the aggregated Dice
$\mathrm{aggDSC} = \tfrac{1}{2}(\mathrm{DSCagg}_{\mathrm{GTVp}} +
\mathrm{DSCagg}_{\mathrm{GTVn}})$ computed after concatenating the
voxels of all test patients per structure; the challenge winner sits
at $0.733$ aggDSC and the widely-quoted nnU-Net reference at $0.727$.

\subsubsection{Secondary longitudinal benchmark}
\label{sec:exp_brats}

\textbf{BraTS-GLI 2024 post-treatment glioma.} The secondary
longitudinal benchmark is BraTS-GLI 2024
post-treatment~\cite{delarosa2024brats}, a brain-tumor follow-up task
with roughly 118 longitudinal pre/post patient pairs and four MR
modalities per visit (t1c, t1n, t2f, t2w). Each voxel carries one of
five labels (background, non-enhancing tumor core, surrounding
FLAIR hyperintensity, enhancing tumor, resection cavity), evaluated as
lesion-wise Dice on the ET, TC, WT, and RC sub-regions. Because
SAM2's image stem accepts three channels, we use the t2f sequence by
default; 4-modality input through an added $1{\times}1$ conv or
sequence-token stacking is left to future work.

\subsubsection{Pretraining corpora (single-timepoint, medical)}
\label{sec:exp_pretrain}

\textbf{OAIZIB-CM knee MRI.} OAIZIB-CM contains 507 DESS knee MR
volumes with five cartilage/bone ROIs (femur, tibia, femoral
cartilage, medial tibial cartilage, lateral tibial cartilage) plus
background~\cite{yao2023cartimorph}. We load it as
\emph{self-paired}: pre $=$ mid, change mask $=$ 0, \texttt{single\_timepoint=True},
which gates off the change-head and consistency losses and exercises
only the Dice term plus the identity prior of the memory attention.

\textbf{MS-Segmentation.} A supplementary 60-volume brain-MRI
multiple-sclerosis lesion parquet set. There is no published
leaderboard on this exact 60-example split, so it is used only as a
brain-MRI intensity-statistics prior: the same self-paired contract as
OAIZIB-CM with \texttt{single\_timepoint=True}, reporting no primary
metric in Section~\ref{sec:results}. Together, the two pretraining
corpora probe whether the frozen SAM2 ViT-H stem transfers to MR
intensity distributions outside head and neck without destabilising
the cross-session module.

\subsubsection{Transfer benchmarks (bitemporal change detection)}
\label{sec:exp_cd}

We transfer the CSM-SAM mechanism, without architectural changes, to
four remote-sensing / disaster bitemporal change-detection benchmarks.

\textbf{LEVIR-CD~\cite{chen2020levircd}.} Binary building change
detection on 0.5 m RGB aerial tiles; 445 / 64 / 128 tile split of
$1024{\times}1024$ images, evaluated on $256{\times}256$ crops with F1
and IoU on the change class.

\textbf{S2Looking~\cite{shen2021s2looking}.} Off-nadir binary building
change detection on rural satellite imagery; 3500 / 500 / 1000 tile
split. Off-nadir viewing geometry and small buildings push SOTA
roughly 20 F1 points below LEVIR-CD, which stresses viewpoint
robustness of the underlying foundation backbone.

\textbf{SECOND~\cite{yang2021second}.} Semantic change detection with
six land-cover classes plus background on $512{\times}512$ tiles;
2968 pairs with a 1:1 train/test split. The evaluation protocol
follows prior work: mean IoU over semantic-change pixels
($\mathrm{mIoU}_{\mathrm{sc}}$) and the Separated Kappa coefficient
(SeK). SECOND is the closest structural analogue of CSM-SAM's output
signature because it requires \emph{both} per-timepoint semantic maps
and a binary change map.

\textbf{xBD / xView2~\cite{gupta2019xbd}.} Building damage assessment
on $1024{\times}1024$ pre- and post-disaster satellite imagery with a
four-level ordinal damage label (no, minor, major, destroyed) plus an
``un-classified'' category. The official xView2 score is
$0.3 \cdot \mathrm{loc\_F1} + 0.7 \cdot \mathrm{dmg\_F1}$, where
$\mathrm{dmg\_F1}$ is a harmonic mean across the four damage classes.

These four datasets together cover binary vs.\ semantic CD, nadir vs.\
off-nadir geometry, per-pixel vs.\ per-building evaluation, and
two-class vs.\ ordinal label sets, which lets us ask whether a single
cross-session memory module survives each of those axis changes without
retraining the SAM2 backbone.

\subsection{Baselines}
\label{sec:exp_baselines}

We compare CSM-SAM against 36 baselines spanning six categories:
naive / sanity, classical medical segmentation, foundation-model
zero-shot, longitudinal / bitemporal medical, binary change detection,
and semantic change detection. The taxonomy follows the landscape
document shipped with the repository (see Section~\ref{sec:related});
each category isolates a different axis along which CSM-SAM is
claimed to differ from prior work.

\textbf{Naive / sanity (5).} \texttt{identity} (pre-RT mask verbatim
as mid-RT prediction), \texttt{zero}, \texttt{random} (Bernoulli at
empirical foreground density), \texttt{copy\_prev\_slice}, and
\texttt{majority\_voxel} (training-set union prior). These establish
the metric floor and quantify how much ``free'' Dice comes from
tumor-free slices or from trivial cross-session reuse; a learned
method that barely beats them has no per-patient signal.

\textbf{Classical medical segmentation (5).} 2D U-Net
(ResNet-34 encoder), DeepLabV3+ (ResNet-50 + ASPP),
Swin UNETR~\cite{hatamizadeh2021swinunetr},
UNETR~\cite{hatamizadeh2022unetr}, and
MedNeXt~\cite{roy2023mednext}. All are strong ImageNet-initialised or
scratch-trained medical specialists with no pre-session signal. They
isolate ``how much does cross-session information buy?'' at a fixed
medical-specialist capacity.

\textbf{Foundation-model zero-shot (6).} SAM~\cite{kirillov2023sam},
SAM2 point prompt, SAM2 video~\cite{ravi2024sam2},
DINOv2~\cite{oquab2024dinov2} with a linear decoder,
CLIPSeg~\cite{lueddecke2022clipseg}, and
TotalSegmentator~\cite{wasserthal2023totalsegmentator}. These freeze a
strong natural-image or medical-image prior and probe whether the
pretraining alone suffices. The delta between them and CSM-SAM
isolates the contribution of the cross-session memory module on top
of the same frozen prior.

\textbf{Longitudinal / bitemporal medical (6).}
MedSAM2~\cite{ma2024medsam2} (within-session memory only),
\texttt{concat\_channels} (pre and mid stacked as a 2-channel input),
\texttt{siamese\_unet}, \texttt{pre\_mask\_prior} (the HNTS-MRG 2024
winner-style pre-RT-mask-as-channel baseline),
\texttt{longisam} (our SAM2-with-6-channel-fusion ablation that
replaces cross-session attention with a $1{\times}1$ conv), and
\texttt{nnunet\_dual\_channel}~\cite{isensee2021nnunet}. Every entry
shares CSM-SAM's goal of using the pre-RT visit but implements that
with simpler mechanisms than learned cross-session token attention;
\texttt{longisam} in particular is the single-axis ablation that
isolates the attention module from the SAM2 backbone.

\textbf{Binary change detection (6).}
FC-Siam-diff and FC-Siam-conc~\cite{daudt2018fcsiam},
SNUNet-CD~\cite{fang2022snunet},
BIT~\cite{chen2021bit},
ChangeFormer~\cite{bandara2022changeformer}, and
TinyCD. These are the dominant remote-sensing binary-CD
architectures: siamese CNNs / transformers that fuse time points by
feature differencing, concatenation, nested skips, channel attention,
or a low-rank transformer-token bottleneck. None of them sits on top
of a 1B-mask-pretrained foundation backbone.

\textbf{Semantic change detection (4).}
Bi-SRNet~\cite{ding2022bisrnet},
SCanNet~\cite{ding2024scannet},
CED-Net / HRSCD~\cite{daudt2019hrscd}, and
Dual-HRNet~\cite{koo2020dualhrnet} on xBD. These predict
\emph{both} per-timepoint semantic maps \emph{and} a change map,
which is structurally the closest analogue to CSM-SAM's
\{mid-mask, change-map\} output; they differ in the backbone (small
from-scratch siamese vs.\ frozen SAM2 ViT-H) and in whether the
cross-time coupling is symmetric (both branches query each other) or
asymmetric (pre writes to a memory bank that mid reads), which is the
axis CSM-SAM explores.

For every baseline above, the source file in \texttt{baselines/}
carries a top-of-file ``uniqueness note'' that states precisely what
CSM-SAM does that the baseline does not --- e.g.\ ``no cross-session
token attention'', ``no auxiliary change head'', ``no weeks-elapsed
embedding''. These notes are summarised in Section~\ref{sec:related}
into the six-axis landscape diagram so that every reported number in
Section~\ref{sec:results} can be mapped back to the specific
architectural axis it ablates.

\subsection{Metrics}
\label{sec:exp_metrics}

\textbf{Medical (primary).} On HNTS-MRG 2024 we follow the challenge
protocol~\cite{wahid2024hntsmrg} and report
$\mathrm{aggDSC} = \tfrac{1}{2}(\mathrm{DSCagg}_{\mathrm{GTVp}} +
\mathrm{DSCagg}_{\mathrm{GTVn}})$, where $\mathrm{DSCagg}$ for each
structure is computed after concatenating all test-set voxels for that
structure (not averaging per-patient Dice). We additionally report
per-structure Dice and the 95th-percentile Hausdorff distance (HD95)
as secondary metrics. On BraTS-GLI 2024 we report lesion-wise Dice on
the four sub-regions (ET, TC, WT, RC) following~\cite{delarosa2024brats}.

\textbf{Binary change detection.} On LEVIR-CD~\cite{chen2020levircd}
and S2Looking~\cite{shen2021s2looking} we report F1 and IoU on the
change class, following prior CD literature so that our numbers are
directly comparable to the baseline table entries
\cite{daudt2018fcsiam,fang2022snunet,chen2021bit,bandara2022changeformer}.

\textbf{Semantic and ordinal change detection.} On
SECOND~\cite{yang2021second} we report the semantic-change mean IoU
($\mathrm{mIoU}_{\mathrm{sc}}$) and the Separated Kappa coefficient
(SeK), matching~\cite{ding2022bisrnet,ding2024scannet}. On
xBD~\cite{gupta2019xbd} we report the official xView2 weighted score
$0.3 \cdot \mathrm{loc\_F1} + 0.7 \cdot \mathrm{dmg\_F1}$, where
$\mathrm{dmg\_F1}$ is itself the harmonic mean over the four damage
classes and buildings in the ``un-classified'' category count toward
localisation but not toward damage.

\textbf{Ablation-specific.} For the change head we report IoU of the
predicted change map against the binary XOR of pre and mid masks. For
the feature-evolution regulariser we report the consistency MSE
between the predicted mid features and the ground-truth SAM2 encoder
features on tumor regions. Both are internal signals, reported only
in Section~\ref{sec:ablations}.

\subsection{Implementation setup}
\label{sec:exp_impl}

\textbf{Inputs.} CSM-SAM operates at $1024 \times 1024$ per slice,
matching SAM2's native resolution~\cite{ravi2024sam2}. 3D MR volumes
are sliced axially and fed to the network one slice at a time during
inference; during training we use both slice-level batches (for the
per-slice Dice and change losses) and sequence-level batches (for the
within-session memory bank that accumulates $M_{\text{mid}}$ across
consecutive mid-RT slices).

\textbf{Parameter budget.} The SAM2 ViT-H image encoder
(${\sim}600$M parameters) is frozen; the mask decoder is frozen except
for its final layer. The trainable stack totals roughly 2M
parameters: \texttt{CrossSessionMemoryAttention} (cross-session
attention to $M_{\text{aug}}$, within-session attention to
$M_{\text{mid}}$, a learned gate), the \texttt{ChangeHead}, the
continuous weeks-elapsed encoder (sinusoidal plus an MLP), and the
final mask-decoder layer.

\textbf{Optimisation.} AdamW with learning rate $10^{-4}$, weight
decay $0.05$, and a cosine schedule with 5-epoch warmup. The combined
loss is $\mathcal{L} = \mathcal{L}_{\text{Dice}} +
0.3 \cdot \mathcal{L}_{\text{CE}}(\text{change\_map}) +
0.2 \cdot \mathcal{L}_{\text{consistency}}$. On single-timepoint
pretraining batches (OAIZIB-CM, MS-Segmentation) the
$\mathcal{L}_{\text{CE}}$ term is gated off by the
\texttt{single\_timepoint} flag, so only the Dice term and the
identity prior of the attention module are exercised.

\textbf{Training schedule.} 200 epochs on a single GPU, with a
slice-level batch size of 4 and a sequence-level batch size of 1. We
run 5-fold cross-validation on HNTS-MRG and report fold-0 held-out
aggDSC in Section~\ref{sec:results}; remaining folds are used for
ablation sensitivity analysis in Section~\ref{sec:ablations}.

\textbf{Test-time augmentation.} At inference we average predictions
over horizontal flip, vertical flip, and $90^{\circ}$ rotation.

\textbf{Synthetic smoke-test data.} For pipeline validation we ship a
15-patient synthetic HNTS-MRG replica, generated by
\texttt{data/download\_hnts\_mrg.py -{}-synthetic}, that reproduces
the voxel spacing, label set, and pre/mid pairing of the real
challenge data without any patient imagery. Whenever real HNTS-MRG
preprocessing is not available on the execution environment, the
hyperparameter sweep in Section~\ref{sec:results} falls back to this
synthetic split; we flag those runs explicitly so they are not
conflated with aggDSC on the true challenge test set.
